\begin{textsample}{NEG Dim 1 – 2012\_c5 – Score 4.00 – 2012\_c5/enade\_2012\_ri\_3badec8761.txt}  \label{ex:f1_neg_321}
CAPÍTULO 5 DISTRIBUIÇÃO DOS CONCEITOS 5.1 PANORAMA NACIONAL DA DISTRIBUIÇÃO DOS CONCEITOS A Tabela 5.1 apresenta a quantidade e distribuição de cursos de Relações Internacionais participantes do ENADE/2012, por faixa de conceito e Grande Região.
A diferença entre os cursos tabulados neste capítulo e no capítulo 2 são os cursos sem conceito, em princípio, aqueles sem alunos concluintes que participassem da prova ou que tivessem tão somente um aluno concluinte11.
Observando -se os dados da Tabela 5.1, nota -se que, dos 78 cursos participantes, 38 ( 48,7 \% ) classificaram -se com conceito 3, o valor modal.
Este foi também o conceito modal em todas as regiões do Brasil ( entre 42,9 \% no Nordeste e 75,0 \% no Centro-Oeste ).
O conceito 4 foi o segundo mais frequente em nível nacional ( 23,1 \%, correspondendo a 18 cursos ) e o conceito 2, o terceiro ( 15,4 \%, correspondendo a 12 cursos ).
Houve, ainda, quatro cursos ( 5,1 \% ) que receberam conceito 5 e outros três cursos que receberam conceito 1 ( 3,8 \% ).
Apenas três dos 78 cursos de Relações Internacionais ficaram sem conceito ( SC ).
A região Norte participou com três cursos ou 3,8 \% do total nacional.
Nesta região, o conceito 3 foi atribuído a dois cursos dos três participantes, o que equivale a 66,7 \% do total regional.
O outro curso participante ficou sem conceito ( SC ).
A região Nordeste participou com sete cursos ou 9,0 \% do total nacional.
Como já comentado, destes, três cursos, 42,9 \% em termos regionais, obtiveram conceito 3, o conceito modal para a região.
O conceito 2 foi atribuído a dois cursos ( 28,6 \% ) e o conceito 4 a um curso ( 14,3 \% ).
Nessa região, somente um dos cursos ficou sem conceito ( SC ). 11 Estes cursos com somente um concluinte foram considerados no capítulo 2. 5.2 CONCEITOS POR CATEGORIA ADMINISTRATIVA E POR GRANDE REGIÃO A Tabela 5.2 apresenta a distribuição dos cursos participantes do ENADE/2012 de Relações Internacionais, por Categoria Administrativa, de acordo com os conceitos por eles alcançados, segundo as Grandes Regiões brasileiras.
Dos 78 cursos participantes, 13 ( 16,7 \% ) eram ministrados em instituições públicas e 65 ( 83,3 \% ), em privadas.
De acordo com as informações da Tabela 5.2, em termos nacionais, entre instituições públicas e privadas : todas as quatro IES que receberam o conceito 5 eram públicas.
Dos 13 cursos participantes de IES públicas, cinco cursos ( 38,5 \% ) receberam o conceito 4, o valor modal.
Nesta categoria, somente um dos cursos ficou sem conceito ( SC ).
Entre os demais cursos de instituições públicas participantes, um obteve conceito 1 ( 7,7 \% da categoria ), dois receberam conceito 3 ( 15,4 \% ) e quatro foram alocados ao conceito 5 ( 30,8 \% ).
Na rede privada, o conceito modal foi 3, com 36 cursos dos 65 da categoria.
Nesta categoria administrativa, dois dos cursos ficaram sem conceito ( SC ).
Entre os demais cursos participantes, dois receberam conceito 1, 12 receberam conceito 2 e o conceito 4 foi atribuído a 13 cursos.
O conceito 5, como já comentado, não foi atribuído a nenhum curso.
Dos 41 cursos participantes da região Sudeste, 18 ( 43,9 \% ) obtiveram conceito 3, o conceito modal.
O conceito 1 foi atribuído a dois cursos ( 4,9 \% ) e o conceito 2, a oito ( 19,5 \% ).
Onze cursos ( 26,8 \% ) receberam o conceito 4 e um outro ( 2,4 \% ) recebeu o conceito 5.
Somente um curso ficou sem conceito ( SC ).
A região Sul contou com cursos distribuídos em todas as faixas de conceitos.
A predominância do conceito 3 foi de 47,4 \%, correspondentes a nove dos 19 cursos participantes na região Sul.
O conceito 2 foi atribuído a dois cursos ( 10,5 \% ) e o conceito 4, a cinco cursos ( 26,3 \% ).
O conceito 1 foi recebido por um único curso ( 5,3 \% ) e o conceito 5, por dois deles ( 10,5 \% ).
Nenhum dos cursos da região Sul ficou sem conceito ( SC ).
Como já comentado, exatamente 3/4 ( 75,0 \%, correspondendo a seis cursos ) dos oito cursos participantes na região Centro-Oeste recebeu conceito 3, o conceito modal.
Os demais cursos foram avaliados com conceitos 4 e 5 ( um curso em \textbf{cada}, 12,5 \% ).
Nesta região nenhum curso ficou sem conceito ( SC ).
Nenhum curso tampouco foi alocado nos conceitos 1 ou 2. 5.2 CONCEITOS POR CATEGORIA ADMINISTRATIVA E POR GRANDE REGIÃO A Tabela 5.2 apresenta a distribuição dos cursos participantes do ENADE/2012 de Relações Internacionais, por Categoria Administrativa, de acordo com os conceitos por eles alcançados, segundo as Grandes Regiões brasileiras.
Dos 78 cursos participantes, 13 ( 16,7 \% ) eram ministrados em instituições públicas e 65 ( 83,3 \% ), em privadas.
De acordo com as informações da Tabela 5.2, em termos nacionais, entre instituições públicas e privadas : todas as quatro IES que receberam o conceito 5 eram públicas.
Dos 13 cursos participantes de IES públicas, cinco cursos ( 38,5 \% ) receberam o conceito 4, o valor modal.
Nesta categoria, somente um dos cursos ficou sem conceito ( SC ).
Entre os demais cursos de instituições públicas participantes, um obteve conceito 1 ( 7,7 \% da categoria ), dois receberam conceito 3 ( 15,4 \% ) e quatro foram alocados ao conceito 5 ( 30,8 \% ).
Na rede privada, o conceito modal foi 3, com 36 cursos dos 65 da categoria.
Nesta categoria administrativa, dois dos cursos ficaram sem conceito ( SC ).
Entre os demais cursos participantes, dois receberam conceito 1, 12 receberam conceito 2 e o conceito 4 foi atribuído a 13 cursos.
O conceito 5, como já comentado, não foi atribuído a nenhum curso.
Na análise por região, observa -se que, na região Norte, as instituições privadas participaram com dois cursos ( 66,7 \% do total regional, a menor proporção entre todas as regiões ), os quais receberam conceito 3.
As instituições públicas participaram com somente um curso ( 33,3 \% do total regional ), que ficou sem conceito ( SC ).
Na região Nordeste, a rede privada concentrou cinco dos sete cursos participantes, equivalentes a 71,4 \% do total da região, dos quais um curso ficou sem conceito ( SC ).
Nesta Categoria Administrativa na região, os demais quatro cursos foram avaliados com conceitos 3 e 4 ( com dois cursos \textbf{cada} um, correspondendo a 40,0 \% ).
As instituições públicas dessa região participaram com dois cursos ( 28,6 \% ), que obtiveram os conceitos 3 e 4 ( sendo um \textbf{cada} ).
Nenhum curso ficou sem conceito ( SC ) nesta combinação de região e categoria.
Na região Sudeste, a proporção de cursos da rede privada, 87,8 \%, foi mais elevada do que nas demais regiões brasileiras, correspondendo a 36 dos 41 cursos participantes.
Nesta categoria, na região Sudeste, o conceito modal foi 3 ( 18 cursos ).
Nesta combinação de Categoria Administrativa e Grande Região, somente um dos cursos ficou sem conceito ( SC ).
Os demais foram avaliados com conceito 1 ( um curso ), conceito 2 ( oito cursos ) e conceito 4 ( outros oito cursos ).
Entre os cinco cursos em instituições públicas na região Sudeste, a categoria modal foi a que recebeu conceito 4 ( três cursos ).
Os demais foram avaliados com os conceitos 1 e 5, com um curso \textbf{cada}.
Na região Sudeste, nesta categoria, nenhum curso ficou sem conceito ( SC ).
As instituições privadas concentraram 16 dos 19 cursos participantes da região Sul, 84,2 \% do total regional.
Desses, oito ficaram com conceito 3, o conceito modal.
Os demais foram avaliados com conceito 1 ( um curso ), conceito 2 ( dois cursos ) e conceito 4 ( cinco cursos ).
Nesta combinação de Categoria Administrativa e Grande Região, nenhum curso ficou sem receber conceito ( SC ).
As instituições públicas na região Sul participaram com três cursos ( 15,8 \% ), dos quais nenhum ficou sem conceito ( SC ).
O conceito modal foi 5, com dois cursos.
Além disso, o conceito 3 foi recebido por um curso.
Na região Centro-Oeste, seis dos oito cursos participantes eram de instituições privadas ( 75,0 \% em termos regionais ), todos avaliados com conceito 3.
Dos dois cursos de instituições públicas, um foi avaliado no conceito 4 e o outro recebeu conceito 5.
Nesta região, nenhum curso ficou sem conceito ( SC ). 5.3 CONCEITOS POR ORGANIZAÇÃO ACADÊMICA E POR GRANDE REGIÃO Na Tabela 5.3 encontra -se a distribuição dos conceitos atribuídos aos cursos participantes do ENADE/2012 na Área de Relações Internacionais, por Organização Acadêmica, segundo as Grandes Regiões brasileiras.
Dos 78 cursos participantes, 37 eram oferecidos em Universidades, 19 em Centros Universitários e os demais 22 em Faculdades.
Esta distribuição corresponde a, respectivamente, 47,4 \%, 24,4 \% e 28,2 \% dos cursos.
De acordo com os dados apresentados, todos os quatro cursos avaliados com conceito 5 eram vinculados a Universidades.
Dois cursos em Universidades ficaram sem conceito ( SC ).
Esse tipo de Organização Acadêmica teve o conceito 3 como modal, com 14 cursos ( 37,8 \% ).
Os demais cursos de Universidades avaliados receberam os conceitos 1 ( dois cursos ), 2 ( quatro cursos ), 4 ( 11 cursos ) e conceito 5 ( outros quatro cursos, como já mencionado ).
Entre os cursos em Centros Universitários, o conceito modal também foi 3, com 12 cursos ( 63,2 \% ).
Neste tipo de organização acadêmica nenhum curso ficou sem conceito ( SC ) ou recebeu o conceito 5.
Os outros cursos neste tipo de Organização Acadêmica receberam os conceitos 1 ( um curso ), 2 ( três cursos ) e 4 ( outros três cursos ).
Nas Faculdades, somente um dos 22 cursos ficou sem conceito ( SC ) e 12 ( 54,5 \% ) receberam o mesmo conceito modal dos demais tipos de Organização Acadêmica, conceito 3.
Dos demais cursos neste tipo de Organização Acadêmica, cinco receberam conceito 2 e quatro cursos, o conceito 4.
Nenhum curso ficou sem conceito ( SC ) ou recebeu os conceitos 1 ou 5.
Considerando -se separadamente as regiões brasileiras, verifica -se que, na região Norte, as Universidades concentraram dois dos três cursos participantes.
Dos cursos em Universidades, um deles ficou sem conceito ( SC ) e o outro recebeu conceito 3.
Os Centros Universitários da região Norte não foram representados por nenhum curso.
As Faculdades participaram com somente um curso, o qual foi avaliado com conceito 3.
Na região Nordeste, as Universidades participaram com quatro dos sete cursos na Área de Relações Internacionais da região.
Um dos cursos em Universidades no Nordeste ficou sem conceito ( SC ), enquanto nenhum recebeu os conceitos 1 e 5.
Os demais receberam conceito 2, 3 e 4, com um curso \textbf{cada}.
Os Centros Universitários contaram com somente um curso participantes na região Nordeste, este avaliado com conceito 3.
As Faculdades foram representadas por dois cursos na região Nordeste, um com conceito 2 e o outro, com conceito 3.
Na região Sudeste, as Universidades concentraram 18 dos 41 cursos da região.
Entre os cursos em Universidades na região, o conceito modal foi 4 com sete cursos, e nenhum curso ficou sem conceito ( SC ).
Os demais cursos receberam os conceitos 1 ( um curso ), 2 ( três cursos ), 3 ( seis cursos ) e 5 ( também um curso ).
Os Centros Universitários participaram com dez cursos na região Sudeste, dos quais cinco obtiveram conceito modal, 3, e nenhum ficou sem conceito ( SC ).
Os demais receberam os conceitos 1 ( um curso ), 2 ( três cursos ) e conceito 4 ( também um curso ).
Nenhum dos cursos nesta combinação de organização acadêmica e região recebeu conceito 5.
As Faculdades foram representadas por 13 cursos na região Sudeste, que se distribuíram nos conceitos 2 ( dois cursos ), 3 ( sete cursos, conceito modal ), 4 ( três cursos ), enquanto um curso ficou sem conceito ( SC ).
Dos 19 cursos da região Sul, nove eram de Universidades, para os quais o conceito modal foi 3, com quatro cursos.
Nesse tipo de organização, nenhum dos cursos ficou sem conceito ( SC ) e os demais receberam os conceitos 1 ( um curso ), 4 ( dois cursos ) e 5 ( outros dois cursos ).
Os Centros Universitários da região Sul tiveram dois dos quatro cursos participantes no conceito 3, e os outros dois cursos receberam conceito 4.
Foram seis os cursos vinculados a Faculdades na região Sul, e nenhum destes ficou sem conceito ( SC ).
Dois desses cursos receberam conceito 2, três cursos receberam conceito 3, o conceito modal, e o outro, conceito 4.
Na região Centro-Oeste, quatro dos oito cursos eram de Universidades.
Nesse tipo de organização, nenhum curso ficou sem conceito ( SC ) e o conceito modal foi 3, com dois cursos.
Os outros dois cursos obtiveram os conceitos 4 e 5, com um curso \textbf{cada}.
Os Centros Universitários da região Centro-Oeste contaram com quatro cursos, todos eles receberam o conceito 3.
As Faculdades na região Centro-Oeste não participaram com nenhum curso.

% matched lemmas: cada
\end{textsample}
